\subsection{Torsion and Annihilators}

Throughout this subsection, $R$ denotes a nonzero commutative ring with identity and $M$ denotes an $R$-module.

\begin{definition} \label{def:module-torsion} \hlt{Torsion}
\begin{enumerate}[label=(\roman*)]
\setlength{\itemsep}{0pt}
\item A nonzero element $v\in M$ is a \hlt{torsion element} if $\exists\,0\neq r\in R$ such that $rv=0$.
\item $M$ is \hlt{torsion-free} if it has no torsion elements.
\item $M$ is a \hlt{torsion module} if every nonzero element of $M$ is torsion.
\item The \hlt{torsion set} of $M$ is $M_{\operatorname{tor}}\defeq\{0\}\cup\{v\in M:v$ is torsion$\}$
\end{enumerate}
\end{definition}

\begin{example} \label{ex:zn-is-torsion-module} \hlt{$\mathbb{Z}/n\mathbb{Z}$ as a Torsion Module}\\
For $n>1$, $\mathbb{Z}/n\mathbb{Z}$ is a $\mathbb{Z}$-module via $k\cdot \overline{a}=\overline{ka}$. For every $\overline{a}\in\mathbb{Z}/n\mathbb{Z}$, $n\overline{a}=\overline{0}$.\\
Hence $\mathbb{Z}/n\mathbb{Z}$ is a torsion module, and $\{\overline{a}\}$ is linearly dependent whenever $\overline{a}\neq\overline{0}$.
\end{example}

\begin{proposition} \label{prop:torsion-submodule-over-domain} \hlt{Torsion Submodule over an Integral Domain}\\
If $R$ is an integral domain, then $M_{\operatorname{tor}}\leq M$, $(M/M_{\operatorname{tor}})_{\operatorname{tor}}=\{0\}$.
\begin{proof}
Let $x,y\in M_{\operatorname{tor}}$ and $r,s\in R$.\\
Choose $a,b\in R\setminus\{0\}$ with $ax=0$ and $by=0$, taking $a=1$ if $x=0$ and $b=1$ if $y=0$.\\
Since $R$ is a domain, $ab\neq0$, and $ab(rx+sy)=rb(ax)+sa(by)=0$. Thus $rx+sy\in M_{\operatorname{tor}}$, so $M_{\operatorname{tor}}\leq M$.\\
Now let $0\neq r\in R$ and suppose $r(m+M_{\operatorname{tor}})=M_{\operatorname{tor}}$. Then $rm\in M_{\operatorname{tor}}$, so $srm=0$ for some $0\neq s\in R$; if $rm=0$, take $s=1$. Since $sr\neq0$, $m\in M_{\operatorname{tor}}$, hence $m+M_{\operatorname{tor}}=M_{\operatorname{tor}}$.\\
\hlt{Structural Commentary}: Domains let annihilators be multiplied without becoming zero.\\
\hlt{Key Insight}: Domains make torsion closed and quotients torsion-free.
\end{proof}
\end{proposition}

\begin{definition} \label{def:module-annihilator} \hlt{Annihilators}
\begin{enumerate}[label=(\roman*)]
\setlength{\itemsep}{0pt}
\item For $v\in M$, the \hlt{annihilator} of $v$ is $\operatorname{Ann}_R(v)\defeq\{r\in R:rv=0\}$.
\item For $N\leq M$, the \hlt{annihilator} of $N$ is $\operatorname{Ann}_R(N)\defeq\{r\in R:rN=\{0\}\}$.
\end{enumerate}
\end{definition}

\begin{proposition} \label{prop:annihilator-basic-properties} \hlt{Annihilator Properties}\\
For $v\in M$ and $N\leq M$:
\begin{enumerate}[label=(\roman*)]
\setlength{\itemsep}{0pt}
\item $\operatorname{Ann}_R(v)\trianglelefteq R$ and $\operatorname{Ann}_R(N)\trianglelefteq R$.
\item $v$ is torsion $\Longleftrightarrow$ $v\neq0$ and $\operatorname{Ann}_R(v)\neq\{0\}$.
\item $A\leq B\leq M \Longrightarrow \operatorname{Ann}_R(B)\subseteq\operatorname{Ann}_R(A)$.
\end{enumerate}
\begin{proof}
\begin{enumerate}[label=\roman*.]
\setlength{\itemsep}{0pt}
\item[$(i)$] Since $0\in\operatorname{Ann}_R(v)$ and, for $a,b\in\operatorname{Ann}_R(v)$ and $r,s\in R$, $(ra+sb)v=r(av)+s(bv)=0$, $\operatorname{Ann}_R(v)\trianglelefteq R$. The same computation for every $n\in N$ gives $\operatorname{Ann}_R(N)\trianglelefteq R$.
\item[$(ii)$]$(\Rightarrow)$ If $v$ is torsion, then $v\neq0$ and $\exists\,0\neq r\in R$ with $rv=0$, so $\operatorname{Ann}_R(v)\neq\{0\}$.
\item[$(ii)$]$(\Leftarrow)$ If $v\neq0$ and $\operatorname{Ann}_R(v)\neq\{0\}$, choose $0\neq r\in\operatorname{Ann}_R(v)$. Then $rv=0$, so $v$ is torsion.
\item[$(iii)$] If $A\leq B\leq M$ and $r\in\operatorname{Ann}_R(B)$, then $rA=\{0\}$; hence $r\in\operatorname{Ann}_R(A)$.
\end{enumerate}
\hlt{Structural Commentary}: Annihilators reverse inclusion.\\
\hlt{Key Insight}: Killing more vectors gives fewer scalars.
\end{proof}
\end{proposition}

\begin{proposition} \label{prop:finitely-generated-torsion-module-annihilator} \hlt{Torsion Generators Give a Nonzero Annihilator over a Domain}\\
Assume $R$ is an integral domain. If $M=\langle\!\langle u_1,\ldots,u_n\rangle\!\rangle_R$, $u_i\in M_{\operatorname{tor}}\ \forall i$, then $\operatorname{Ann}_R(M)\neq\{0\}$.
\begin{proof}
Choose $0\neq a_i\in R$ with $a_iu_i=0$, and set $a\defeq\prod\limits_{i=1}^{n}a_i$. Since $R$ is a domain, $a\neq0$.\\
For each $i$, $au_i=0$, hence $aM=\{0\}$ because the $u_i$ generate $M$. Thus $0\neq a\in\operatorname{Ann}_R(M)$.\\
\hlt{Structural Commentary}: Finite generation lets many annihilators be compressed into one product.\\
\hlt{Key Insight}: A finite torsion generating set gives a global order for the module.
\end{proof}
\end{proposition}
